% Fragment: §1-§2 of Bounded Deontics. Preamble, macros, theorem envs live in bd-draft.tex.
% (For standalone compilation, prepend the bd-draft.tex preamble.)

\section{Introduction}
\label{sec:intro}

Every legal ``must'' invites a retort: \emph{or else what?} A statute or a contract
tells us we \emph{must} file within thirty days, \emph{must} return the book, \emph{must}
give notice; the vernacular question asks what stands behind the word---what consequence,
and in service of what end. That question is not impertinent. It is the one a court asks
when the relationship sours, and it is the one a programmer asks when a specification says
a system ``must'' respond. This paper takes the question seriously and follows it to a
formal account of obligation in which a ``must'' is \emph{bounded} in two senses we keep
apart: bounded in \emph{domain}---it is true only relative to a frame, a goal, and a
repertoire of available acts---and, within that domain, \emph{binding} only to a bounded
degree, a force that can be measured and can fall to zero. \emph{Bounded}, not \emph{bound}:
the modal is hemmed in before it grips, and the grip, when it comes, is itself finite.

The starting distinction is between the \emph{letter} and the \emph{spirit} of a rule. A
program executes; a specification judges. The letter is the executable artefact---the
states a contract can be in and the actions it admits; the spirit is the property that
artefact is supposed to satisfy. In model checking these inhabit different semantic
levels: the implementation is a transition system, and a temporal-logic formula is not a
rival implementation but a statement \emph{about} it. Law has the same two levels. The
letter is the rule as written; the spirit is the purpose any compliant procedure must
serve. The spirit is what survives refactoring: one may amend the forms, the deadlines,
the administering agency, and still mean to preserve the same underlying objective.

This two-level reading has a venerable jurisprudential shadow. For Austin and for Holmes,
law is best seen through the eyes of the \emph{bad man}, who ``cares only for the material
consequences'' and treats a duty as a prediction of sanction \src{Holmes 1897}. For Hart,
that view misses what he called the \emph{internal point of view}: subjects of a
functioning legal system do not merely predict the application of rules but \emph{use}
them as standards for appraising conduct \src{Hart 1961}. The contrast has acquired fresh
urgency: as autonomous AI agents become subjects, users, and even producers of law, they
threaten to be the perfect Holmesian bad man---optimisers that respect the letter while
gaming the spirit---which is precisely why we need formal tools for the spirit and not
only the letter \src{Kolt 2026}.

What goes wrong when the spirit is squeezed out? Two illustrations, one empirical and one
literary, fix the phenomenon. Gneezy and Rustichini report that a day-care centre that
imposed a \emph{fine} for late pickup saw lateness \emph{rise}: parents read the fine not
as a sanction reinforcing a duty but as a price for a service, and ``a fine is a price''
\src{Gneezy \& Rustichini 2000}. Compiling an obligation into a priced transition can
dissolve the obligation. The literary case runs the compiler the other way. In Saki's
\emph{The Schartz-Metterklume Method}, told ``You must be Miss Hope, the governess,'' Lady
Carlotta replies ``Very well, if I must I must''---deliberately hearing an \emph{epistemic}
``must'' (you surely are) as a \emph{deontic} one (you are required to be), and assuming
the role. The same modal word carries different forces; which force is in play depends on
a goal the speaker leaves implicit. Both cases teach the same lesson: a ``must'' is bounded
by a goal---and, \emph{within} those bounds, its binding force is \emph{agent-relative}.
(Section~\ref{sec:boundary} pushes this to its limit with a protester for whom the
threatened sanction is the goal.)

Consider, finally, a case with no rule in it at all. A restaurant patron \emph{may} raise a
hand, fold a napkin, or let a spoon fall---all equally permitted, none prescribed; the
``letter'' here is a bare menu of available acts, with no obligation written into any of
them. Yet if he wants another glass of wine, and the convention of the house is that wine is
summoned from the waiter and not fetched by the patron from the cellar, then he \emph{must}
raise his hand. No clause obliges him: the ``must'' is manufactured entirely by a goal
together with the absence of any other route to it, and it singles out one permitted act
from a sea of equally permitted others by that act's \emph{relation} to the goal alone. We
take this derived, rule-free ``must'' as our paradigm---obligation read \emph{off} the
letter, not written \emph{into} it (Section~\ref{sec:nec}).

\paragraph{Thesis.} Deontic force is not a primitive but a quantity \emph{derived} from the
letter and \emph{bounded}---in the domain where it is defined, and in the magnitude to which
it binds. On the account we develop: (i) the goal-relativised
``must'' is a relation \emph{derived} from the letter---a necessity relation over a
transition system---rather than an operator annotated onto an obligation; (ii) layering a
party-indexed ordering source (the agent's goals) over that necessity yields a
teleological \MUST{}; (iii) two distinguished orderings---the legal system's and the
agent's---generate, by their \emph{divergence}, both Holmes's bad man and Hart's internal
point of view; and (iv) the \emph{force} of a ``must'' is the gap, in the agent's
ordering, between complying and breaching---which is why a small fine, or a desired jail
cell, can drain or even reverse it.

\paragraph{Contributions, and an honest posture.} We claim no new logical primitive. Every
ingredient below is established prior art, and we concede each by name in
Section~\ref{sec:background}: the reduction of obligation to a modality plus a violation
\src{Anderson 1958; Meyer 1988}; necessity-to-a-goal as a graph dominator, known in
planning as a \emph{landmark} \src{Hoffmann, Porteous \& Sebastia 2004}; betterness-and-
priority orderings \src{van Benthem, Grossi \& Liu 2010}; and obligation as strategic
ability \src{Wooldridge \& van der Hoek 2005}. Our contribution is the \emph{composition}
and its \emph{jurisprudential reading}: (a)~the Holmes-and-Hart unification on a single
ordering divergence; (b)~a party-indexed ordering composed specifically with a
\emph{dominator} necessity; (c)~the sanction/goal duality as the \emph{sign} of an outcome
under the agent's ordering; (d)~the force of an obligation as an ordinal preference gap;
and (e)~an executable instantiation of all of this in the legal DSL L4. The reduction is
old; the synthesis, and the legal payoff, are the work.

\paragraph{Roadmap.} Section~\ref{sec:background} places the account in its three
literatures. Section~\ref{sec:nec} defines the derived necessity relation---the bounded
\MUST{} as goal-domination---and is deliberately party-neutral. Section~\ref{sec:boundary}
adds the party-indexed ordering, derives the Holmes/Hart divergence and the force-gap, and
maps the boundary within which the whole reduction is faithful.
Section~\ref{sec:tooling} reports what of this is realised in the L4 toolchain today and
what remains; Section~\ref{sec:eval} works one example---perfection of a security interest
under UCC Article~9---end to end; and Section~\ref{sec:conclusion} concludes.

\section{Background: a reduction with a long genealogy}
\label{sec:background}

The idea that an obligation can be reduced to other, better-understood notions---a
modality, a goal, a sanction, an ordering---is not ours; it is one of the oldest moves in
the formal study of norms. We rehearse it deliberately, because the contribution of this
paper is to compose these strands, not to replace them, and an honest synthesis must first
own its debts.

\subsection{Reducing ``ought'' to a modality and a violation}
The deflationary tradition in deontic logic defines obligation in terms of necessity plus
a marker of breach. Anderson reduces ``it ought to be that $\varphi$'' to ``necessarily,
$\neg\varphi$ brings about the sanction''---$O\varphi \leftrightarrow \Box(\neg\varphi \to
S)$ for a propositional sanction constant $S$ \src{Anderson 1958}. This is the
``or else what?'' of our introduction, formalised in 1958. Meyer pushes the same intuition
into a \emph{state machine}: in dynamic deontic logic an action is forbidden iff every way
of performing it leads to a distinguished violation state, $F\alpha \leftrightarrow
[\alpha]V$ \src{Meyer 1988}. The violation atom $V$ is exactly the neutral
\textsc{breach} terminal our object level will use in Section~\ref{sec:nec}: a marker, not
a verdict, whose moral weight is supplied by context. We descend from this line directly,
and say so.

\subsection{Goal-relative modality}
A parallel tradition in linguistic semantics holds that a modal is never an unindexed
primitive. On Kratzer's standard account a ``must'' is necessity relative to a
\emph{modal base} (the relevant circumstances) and an \emph{ordering source} (a set of
goals or norms) \src{Kratzer 1991}. The point is sharpened by \emph{anankastic
conditionals}---``if you want to go to Harlem, you must take the A train''---in which the
agent's stated goal openly supplies the ordering source \src{von Fintel \& Iatridou 2005;
Condoravdi \& Lauer 2016}. The same move appears in recent logics of \emph{instrumental}
obligation, where the obligation takes the goal it serves as an argument
\src{Yan \& He 2025}, and in accounts of obligation as optimal goal satisfaction, where a
sanction is the worse disjunct of a goal \src{Kowalski \& Satoh 2018}. At the limit of this
strand the modal ceases to be a necessity at all and becomes an \emph{optimum}: a
conditional ought is read as the act that maximises a (qualitatively ranked) expected
utility \src{Pearl 1993}. The two reductions of this section thus mark the poles of a single
spectrum---from Anderson's box over \emph{every} path, through Kratzer's box over the
\emph{best} worlds, to Pearl's pure best-choice---along which one slides by exposing first an
ordering and then a probability. We adopt the
ordering-source idea wholesale; our departure (Section~\ref{sec:ordering}) is to index it
by \emph{party} and to compose it with a graph relation rather than a best-worlds selection,
and to sit at the \emph{necessity} pole---an ordinal ordering and \emph{no} probability---for
a reason the tooling section makes good (Section~\ref{sec:tooling}).

\subsection{The jurisprudence of sanction and acceptance}
The legal-theoretic counterpart is the long argument between sanction-based and
acceptance-based accounts of legal obligation. Austin's command theory makes law a
sovereign's commands backed by the threat of an ``evil''; Holmes's bad man is its
behavioural shadow \src{Austin 1832; Holmes 1897}. Hart's decisive reply distinguishes
\emph{being obliged} (the gunman's victim) from \emph{having an obligation}, and locates
the latter in the internal point of view \src{Hart 1961}. Silk gives this contrast a
formal semantics, reading Hart's internal/external distinction as a choice of which
ordering-source variable a modal is evaluated against \src{Silk 2019}; Kolt transposes the
whole debate to autonomous AI agents and ties the internal point of view to respect for a
law's spirit over its letter \src{Kolt 2026}. On the empirical side, the behavioural
economics of \emph{motivation crowding} and the expressive theory of law confirm that an
obligation and a price are not interchangeable: making the consequence explicit can
\emph{weaken} the norm \src{Gneezy \& Rustichini 2000; Frey \& Jegen 2001; Cooter 1984}.

\subsection{Necessity as domination, and obligation as strategic ability}
Finally, two formal toolkits supply the machinery we compose. First, ``what must happen on
the way to a goal'' is, in automated planning, the theory of \emph{landmarks}: a fact or
action that lies on every plan reaching the goal \src{Hoffmann, Porteous \& Sebastia 2004},
computable via graph \emph{dominators} \src{Lengauer \& Tarjan 1979}. Second, when goals
are pursued among strategic parties, ``what an agent can bring about'' is the province of
alternating-time temporal logic \src{Alur, Henzinger \& Kupferman 2002}, in which
obligation has been recast as a coalitional \emph{ability} \src{Wooldridge \& van der Hoek
2005}. To these we add the betterness-plus-priority view of deontics as an ordering
problem \src{Hansson 1969; van Benthem, Grossi \& Liu 2010}.

\subsection{Hohfeld and the theory of normative positions}
\label{sec:normative-positions}
The strands above concern conduct \emph{within} a fixed rulebook. But law also empowers
certain actors to \emph{change} the rulebook---a legislature enacts, a judge sentences, a
party to a contract waives or assigns---and this second order has a genealogy of its own on
which the present paper leans. Its vocabulary is Hohfeld's: his two Yale articles
\src{Hohfeld 1913; Hohfeld 1917} decompose the loose word ``right'' into eight conceptions,
arranged by two relations---\emph{jural correlatives} (one relation seen from the two
parties' sides: Alice's \emph{claim} that Bob mow the lawn is Bob's \emph{duty} to mow it)
and \emph{jural opposites} (the presence of one is the absence of the other for a party: a
\emph{privilege} to enter negates a \emph{duty} not to enter). The eight fall into two
groups. The \emph{first}---claim, duty, privilege, no-right---concerns what a party may or
must \emph{do}, and maps onto the party-indexed \MUST{}/\MAY{} of this paper. The
\emph{second}---power, liability, immunity, disability---concerns the capacity to
\emph{alter} a first-group position: a \emph{power} changes another party's position,
\emph{liability} is susceptibility to such change, and \emph{immunity} and \emph{disability}
are their opposites.

That the second group is ``one order up'' from the first---operating \emph{on} legal
relations rather than \emph{within} them---is the received view, not a discovery: it runs
from Fitch's revision of Hohfeld \src{Fitch 1967} and Makinson's remarks on the
Kanger--Lindahl formalisation \src{Makinson 1986} to recent surveys that call the second
group the class of \emph{higher-order rights} \src{Markovich 2020}. The formal lineage is
the \emph{theory of normative positions} founded by Kanger and expounded by Lindahl
\src{Kanger 1971; Lindahl 1977; Sergot 2013}, which composes a deontic operator with a
``sees to it that'' action modality to enumerate the logically possible positions---Hohfeld's
first group among them. Tellingly, that tradition \emph{saw} the second group but set it
aside: power concerns ``changes of legal relations,'' listed as a well-documented limitation
to be supplied later \src{Sergot 2001}. Jones and Sergot's \emph{counts-as} conditional
later gave power its standard logical account, on which a designated act \emph{counts as} the
creation or modification of an institutional relation and power is emphatically not reducible
to permission \src{Jones \& Sergot 1996}. We take this normative-positions lineage as shared
background; its development into a typed, position-changing \emph{operator}, and our narrow
contribution there, are deferred to Section~\ref{sec:hohfeld-higher-order}.

\subsection{Position}
Each strand supplies a single ingredient---but singly, and in a form that leaves the
others out. The deontic reductions give necessity-plus-violation but no goal-ordering; the
linguistic and instrumental accounts give a goal-indexed modality but attach the goal
\emph{to} the obligation and over a best-worlds, not a graph, semantics; the jurisprudence
names the Holmes/Hart divide but does not compute with it; the planning and strategic
toolkits give exactly the right notion of necessity but no normative reading. No strand
composes a \emph{party-indexed ordering} over a \emph{graph-dominator necessity}, and none
turns the divergence between the law's ordering and the agent's into the formal locus of
the bad-man/internal-point-of-view debate. That composition is the subject of
Sections~\ref{sec:nec}--\ref{sec:boundary}; its executable form in L4 is the subject of
Section~\ref{sec:tooling}.
